%! Author = sriadityavedantam
%! Date = 3/25/26

\section{The Limit of a Sequence}
\lesson{6}{Monday 25 August 2025 9:10}{Day 6}

\begin{definition}[Sequence]
    A sequence on $\rr$ is a function $\n\to\rr$.
\end{definition}

\begin{definition}[Sequence Convergence]
    A sequence $(x_n)$ converges to $L\in\rr$ if for every $\epsilon >0$ there exists $N\in\n$ such that for all $n > N$,
    \[
        \abs{x_n - L} < \epsilon.
    \]
\end{definition}

\begin{corollary}{
    This is true if and only if for every $m\in\n$, there exists $N\in\n$ such that for all $n > N$,
    \[
        \abs{x_n- L} < \dfrac{1}{m}.
    \]
}
    From the original definition of sequence convergence, one direction is immediate.
    For the converse, let $\epsilon > 0$.
    By the Archimedian property, there exists $m\in\n$ such that $\dfrac{1}{m} < \epsilon$.
    Then for sufficiently large $n$, $\abs{x_n - L} < \dfrac{1}{m} < \epsilon$
\end{corollary}

\setcounter{section}{2}
\section{The Monotone Convergence Theorem and a First Look at Infinite Series}

\begin{definition}[Monotonically Increasing/Decreasing]
    A sequence $(x_n)$ is monotonically increasing if:
    \[
        x_1\leq x_2\leq x_3\leq \ldots
    \]
    Or $x_n\leq x_{n+1}$ for all $n\in\n$.
    Similarly, for monotonically decreasing if $x_n\geq x_{n+1}$.
    A function is monotonic if any of these hold.
\end{definition}

\begin{definition}[Bounded]
    $(x_n)$ is bounded if there exists $C\in\rr$ such that $\abs{x_n}\leq C$ for all $n\in\n$.
\end{definition}

\begin{theorem}[Monotonic Convergence Theorem]{
    A bounded monotonic sequence converges.
}
    Assume without loss of generality that $(x_n)$ is increasing.
    Let $S\coloneqq\{x_n : n\in\n\}$.
    Then $S$ is bounded above.
    Let $L\coloneqq\sup(S)$.
    We show $(x_n)\to L$.
    Let $\epsilon > 0$.
    Then there exists $x_N\in S$ such that $x_N > L - \epsilon$, otherwise $L-\epsilon$ is an upper bound for $S$.
    For all $n > N$, since $(x_n)$ is increasing, $x_n\geq x_N > L - \epsilon$.
    Also $x_n\leq L$ for all $n$, so
    \[
        \abs{L - x_n} = L - x_n < \epsilon.
    \]
\end{theorem}

\begin{proposition}{
    If $(x_n)\to L$, then $(-x_n)\to -L$.
}
    Given $\epsilon > 0$, there exists $N\in\n$ such that for all $n > N$, $\abs{x_n - L} < \epsilon$.
    Then for all $n > N$,
    \[
        \abs{(-x_n) - (-L)} = \abs{x_n - L} < \epsilon.
    \]
\end{proposition}

\begin{proposition}{
    If $(x_n)\to L$, and $(x_n)\to M$, then $L = M$.
}

    Given $\epsilon > 0$, there exist $N_1,N_2\in\n$ such that for all $n > N_1$, $\abs{x_n - L} < \dfrac{\eps}{2}$ and for all $n > N_2$, $\abs{x_n - M} < \dfrac{\eps}{2}$.
    Set $N\coloneqq\max\{N_1,N_2\}$.
    Then for $n >N$:
    \begin{align*}
        \abs{L - M} &= \abs{(L - x_n) + (x_n - M)}\\
        &\leq \abs{L - x_n} + \abs{x_n - M}\\
        &< \dfrac{\eps}{2} + \dfrac{\eps}{2}\\
        &= \eps.
    \end{align*}
    Since $\eps$ is arbitrary, $L = M$.
\end{proposition}

\begin{proposition}{
    If $(x_n)\to L$ and $(y_n)\to M$, then $(x_n + y_n)\to L + M$.
}
    Let $N_1,N_2\in\n$ such that for all $n > N_1$, $\abs{x_n - L} < \dfrac{\eps}{2}$ and for all $n > N_2$, $\abs{y_n - M} < \dfrac{\eps}{2}$.
    Set $N\coloneqq\max\{N_1,N_2\}$.
    Then for all $n>N$:
    \begin{align*}
        \abs{(x_n + y_n) - (L + M)} &= \abs{(x_n - L) + (y_n - M)}\\
        &\leq \abs{x_n - L} + \abs{y_n - M}\\
        &< \dfrac{\eps}{2} + \dfrac{\eps}{2}\\
        &= \eps.
    \end{align*}
\end{proposition}

\begin{definition}[Triangle Inequality]
    \begin{align*}
        \abs{a + b} &\leq \abs{a} + \abs{b}\\
        \abs{\abs{a} - \abs{b}} &\leq \abs{a - b}.
    \end{align*}
\end{definition}

\begin{proposition}{
    If $x_n\to L$, then $(\abs{x_n})\to \abs{L}$.
}
    Let $\epsilon > 0$.
    Then there exists $N\in\n$ such that for all $n > N$, $\abs{x_n - L} < \epsilon$.
    For all $n > N$:
    \begin{align*}
        \abs{\abs{x_n} - \abs{L}} &\leq \abs{x_n - L}\\
        &< \epsilon.
    \end{align*}
\end{proposition}